% ============================================================
% Chapter 5: Development
%   Section: News Acquisition: Discovery and Extraction
%     Subsection: Metadata and Publication Time
% Included from chapters/05-development/02-news-acquisition/news-acquisition.tex
% ============================================================

\subsection{Metadata and Publication Time}
\label{subsec:scraping-metadata}

\paragraph{Metadata from the same page.} A strategy can find the text
of an article and still miss its title, author, date, summary or lead
image. When that happens, the missing fields are read from the same
HTML with BeautifulSoup, JSON-LD first and then meta tags, at no extra
request, and never overwriting what the winning strategy found. The
rule came from an incident: trafilatura~2 leaves the title, author and
date out unless asked for its metadata, and the strategy's tests
mocked the call, so every article in the data lake had been analysed
with an empty title, and every retrieval fallback built on the headline
was dead. One test now runs the extractor for real. On the probe of 25~September~2026, 53 of the 60
sampled articles extracted, all of them with a title, 52 with a date
and 44 with an author.

\paragraph{Publication time.} Both extractors reduce the publication
time to a date, and discovery originally threw away the feed's time, so
the delay between publication and reception could only be measured in
days. Since 4~October~2026 three times are kept:

\begin{itemize}
    \item \textbf{The feed's time} for each discovered URL
    (Section~\ref{subsec:scraping-discovery}).
    \item \textbf{The page's own time}, to the minute, when the page
    states one (\path{article:published_time}, the JSON-LD
    \texttt{datePublished}, or a \texttt{<time datetime>} element). A
    bare date is not a time, and neither is exactly midnight, which is
    how date-only sites write it. This costs one more parse of the page
    when trafilatura won, so it is done only for articles that are
    stored, not for the dozens of evidence pages a claim reads. Six of
    seven live articles sampled that day stated their time to the
    second.
    \item \textbf{When ingestion first saw the URL}, recorded by host
    and path in the lake's statistics (\path{lake/stats/sightings.json}),
    deferred URLs included.
\end{itemize}

Section~\ref{subsec:scraping-observability} describes the report built
from them.
